% GENERATED FILE. Edit canonical lessons, then run npm run generate:books.
% canonical-source-hash: fnv1a64:35d1142ca72bd292
% canonical-lessons: GU-C105-sadhan, GU-C105-rakhvu, GU-C105-dharavvu, GU-C105-jitvu, GU-C105-harvu

\chapter{Tool, To keep, To possess, To win, To lose}
\label{ch:gu-tool-to-keep-to-possess-to-win-to-lose}
\hlchaptermodality{105}

\begin{chapteropening}
\textbf{By the end of this chapter:} \emph{I can say a tool, to keep, to possess, to win and to lose.}

Complete the last lesson of chapter 105: i can say a tool, to keep, to possess, to win and to lose.
\end{chapteropening}

\section[sādhan]{\gu{સાધન} (sādhan) — a tool, a means}

\label{lesson:GU-C105-sadhan}



\begin{quote}
Before the new one: say the Gujarati for a qualification, then the Gujarati for an art.
\end{quote}

\subsection*{You'll want to know: \gu{સાધન}}
\textbf{\gu{સાધન}} (\emph{sādhan}) — \textquotedblleft{}a tool, a means\textquotedblright{}.

\begin{grammarlens}[title={What you've built}]
One more word for this chapter.
\end{grammarlens}

\subsection*{Your turn}
\begin{itemize}
\raggedright
  \item \emph{Say it:} \emph{sādhan}
  \item \emph{Say it:} \emph{sādhan}, once more
  \item \emph{Say it:} say \emph{sādhan}
  \item \emph{Recall:} say the Gujarati for a qualification, then the Gujarati for an art, then say \emph{sādhan} again
\end{itemize}

\begin{tcolorbox}[breakable,colback=teal!4,colframe=teal!35!black,title={Before you move on}]
What does \gu{સાધન} mean? (\textquotedblleft{}A tool, a means\textquotedblright{}.) Say it once more.
\end{tcolorbox}
\section[rākhvũ]{\gu{રાખવું} (rākhvũ) — to keep}

\label{lesson:GU-C105-rakhvu}



\begin{quote}
Before the new one: say the Gujarati for an art, then the Gujarati for a tool.
\end{quote}

\subsection*{You'll want to know: \gu{રાખવું}}
\textbf{\gu{રાખવું}} (\emph{rākhvũ}) — \textquotedblleft{}to keep\textquotedblright{}.

\begin{grammarlens}[title={What you've built}]
One more word for this chapter.
\end{grammarlens}

\subsection*{Your turn}
\begin{itemize}
\raggedright
  \item \emph{Say it:} \emph{rākhvũ}
  \item \emph{Say it:} \emph{rākhvũ}, once more
  \item \emph{Say it:} say \emph{rākhvũ}
  \item \emph{Recall:} say the Gujarati for an art, then the Gujarati for a tool, then say \emph{rākhvũ} again
\end{itemize}

\begin{tcolorbox}[breakable,colback=teal!4,colframe=teal!35!black,title={Before you move on}]
What does \gu{રાખવું} mean? (\textquotedblleft{}To keep\textquotedblright{}.) Say it once more.
\end{tcolorbox}
\section[dharāvvũ]{\gu{ધરાવવું} (dharāvvũ) — to possess, to hold}

\label{lesson:GU-C105-dharavvu}



\begin{quote}
Before the new one: say the Gujarati for a tool, then the Gujarati for to keep.
\end{quote}

\subsection*{You'll want to know: \gu{ધરાવવું}}
\textbf{\gu{ધરાવવું}} (\emph{dharāvvũ}) — \textquotedblleft{}to possess, to hold\textquotedblright{}.

\begin{grammarlens}[title={What you've built}]
One more word for this chapter.
\end{grammarlens}

\subsection*{Your turn}
\begin{itemize}
\raggedright
  \item \emph{Say it:} \emph{dharāvvũ}
  \item \emph{Say it:} \emph{dharāvvũ}, once more
  \item \emph{Say it:} say \emph{dharāvvũ}
  \item \emph{Recall:} say the Gujarati for a tool, then the Gujarati for to keep, then say \emph{dharāvvũ} again
\end{itemize}

\begin{tcolorbox}[breakable,colback=teal!4,colframe=teal!35!black,title={Before you move on}]
What does \gu{ધરાવવું} mean? (\textquotedblleft{}To possess, to hold\textquotedblright{}.) Say it once more.
\end{tcolorbox}
\section[jītvũ]{\gu{જીતવું} (jītvũ) — to win}

\label{lesson:GU-C105-jitvu}



\begin{quote}
Before the new one: say the Gujarati for to keep, then the Gujarati for to possess.
\end{quote}

\subsection*{You'll want to know: \gu{જીતવું}}
\textbf{\gu{જીતવું}} (\emph{jītvũ}) — \textquotedblleft{}to win\textquotedblright{}.

\begin{grammarlens}[title={What you've built}]
One more word for this chapter.
\end{grammarlens}

\subsection*{Your turn}
\begin{itemize}
\raggedright
  \item \emph{Say it:} \emph{jītvũ}
  \item \emph{Say it:} \emph{jītvũ}, once more
  \item \emph{Say it:} say \emph{jītvũ}
  \item \emph{Recall:} say the Gujarati for to keep, then the Gujarati for to possess, then say \emph{jītvũ} again
\end{itemize}

\begin{tcolorbox}[breakable,colback=teal!4,colframe=teal!35!black,title={Before you move on}]
What does \gu{જીતવું} mean? (\textquotedblleft{}To win\textquotedblright{}.) Say it once more.
\end{tcolorbox}
\section[hārvũ]{\gu{હારવું} (hārvũ) — to lose (a game)}

\label{lesson:GU-C105-harvu}



\begin{quote}
Before the new one: say the Gujarati for to possess, then the Gujarati for to win.
\end{quote}

\subsection*{You'll want to know: \gu{હારવું}}
\textbf{\gu{હારવું}} (\emph{hārvũ}) — \textquotedblleft{}to lose (a game)\textquotedblright{}.

\begin{grammarlens}[title={What you've built}]
One more word for this chapter.
\end{grammarlens}

\subsection*{Your turn}
\begin{itemize}
\raggedright
  \item \emph{Say it:} \emph{hārvũ}
  \item \emph{Say it:} \emph{hārvũ}, once more
  \item \emph{Say it:} say \emph{hārvũ}
  \item \emph{Recall:} say the Gujarati for to possess, then the Gujarati for to win, then say \emph{hārvũ} again
\end{itemize}

\begin{tcolorbox}[breakable,colback=teal!4,colframe=teal!35!black,title={Before you move on}]
What does \gu{હારવું} mean? (\textquotedblleft{}To lose (a game)\textquotedblright{}.) Say it once more.
\end{tcolorbox}
